problem:  
Let \(X\) be a smooth Fano manifold and \(\omega_t\) the normalized Kähler–Ricci flow  
\[
\frac{\partial \omega_t}{\partial t} = -\operatorname{Ric}(\omega_t) + \omega_t, \qquad \omega_{t=0} = \omega_0 \in c_1(X),
\]
with potentials \(\omega_t = \omega_0 + i\partial\bar{\partial}\varphi_t\) normalized by \(\int_X e^{-\varphi_t}\omega_0^n = \int_X \omega_0^n\).  
The functionals \(\mathcal{M}(\varphi_t)\) (Mabuchi energy) and \(\mathcal{D}(\varphi_t)\) (Ding functional) both evolve along the flow as nonincreasing convex functions of time.  

Suppose \(\|\varphi_t\|_{C^0}\) remains uniformly bounded.  One may then consider rescaled paths  
\[
\psi_T(s) := \varphi_{T + s} - \varphi_T + c_T, \qquad s\in[0,1],
\]
normalized so that \(\int_X e^{-\psi_T(0)}\omega_0^n = \int_X \omega_0^n\) and \(c_T\) is a constant making the mean zero.  Assume that after taking a subsequence, the curves \(\psi_T(s)\) converge in the weak \(L^1\)-topology of \(\omega_0\)-plurisubharmonic functions to a convex geodesic ray \(\psi_\infty(s)\) in the space \(\mathcal{E}^1(X,\omega_0)\).  

**Question:**  
Is the non–Archimedean metric associated to the ray \(\psi_\infty\) (in the sense of Boucksom–Hisamoto–Jonsson) always the one that *maximizes* the non–Archimedean slope of the Ding functional, i.e. corresponds to the *maximally destabilizing filtration* of \(X\)?  Equivalently, does the Kähler–Ricci flow asymptotically “select” a canonical geodesic ray whose non–Archimedean limit realizes the maximal instability direction of \(X\) among all degenerations?

Why is it a "good" problem:  
This reformulation is a natural and mathematically precise question at the interface of Kähler–Ricci flow dynamics, pluripotential theory, and K‑stability.  It avoids ill‑defined probabilistic notions by focusing on weak \(L^1\) limits of translated potential paths, which are rigorously meaningful in the metric geometry of the space \(\mathcal{E}^1(X,\omega_0)\).  The problem proposes a new direction: connecting the *asymptotic geodesic slope* extracted from long‑time Kähler–Ricci flow behavior with the *non‑Archimedean geometric data* of filtrations realizing maximal Ding energy descent.

If true, this would yield a canonical analytic–geometric process selecting the “most unstable degeneration” of a Fano manifold directly from its Ricci flow evolution, thus forging a deep bridge between parabolic PDE theory and non‑Archimedean algebraic geometry.  The statement is simple to express but encodes several foundational theories—convexity of energy functionals, geodesic geometry of Kähler metrics, and valuation‑theoretic K‑instability—making it a profoundly rich and forward‑looking research problem.